\section{Complete Hermite flags and finite spectral scales}\label{sec:classical-flag}
We separate two classical coordinate facts from the attainable posterior
argument. Throughout the confluent pairing, the scalar parameter is in
$[0,1]$ and $\mu$ is a fixed full-support probability measure. No density
of this measure is required.

\begin{lemma}[Arbitrary-prefix Newton--Hermite principle]\label{lem:hermite-prefix}
Let $y_1,\ldots,y_p$ be any ordered multiset of nodes, with repetitions
not necessarily adjacent. The functionals
$\lambda_jf=[y_1,\ldots,y_j]f$, $1\le j\le p$, form a basis of
the Hermite data at the full multiset: at each distinct node $y$, all
derivatives of orders $0,\ldots,m_y-1$ are included, where $m_y$ is
its multiplicity. Consequently, for $H>0$, the functions
$[y_1,\ldots,y_j](z\mapsto t^{Hz})$ span exactly
\[
 \spanop\{t^{Hy}(\log t)^k:
                 y\text{ distinct},\ 0\le k<m_y\}.
\]
\end{lemma}
\begin{proof}
Every prefix divided difference is determined by the complete Hermite
data at that prefix, so it belongs to the full Hermite data space.
The latter has dimension $p$: a polynomial of degree less than $p$
with all these data zero is divisible by
$\prod_y(z-y)^{m_y}$ and hence is zero. On the polynomials
$1,z,\ldots,z^{p-1}$, $\lambda_j$ annihilates degrees below $j-1$
and takes value one on $z^{j-1}$. Its evaluation matrix is therefore
upper triangular with diagonal one, also at repeated nodes by the Hermite
extension. The $p$ functionals are independent and span the data space.
Substitution of $t^{Hz}$ turns the $k$th derivative at $y$ into
$H^kt^{Hy}(\log t)^k$, with a nonzero constant factor. This proves
the asserted function span.
\end{proof}

This is the prefix interpolation principle of
\cite[Proposition 7, p.~48]{deBoor}. In particular, an arbitrary initial
multiset gives complete multiplicity blocks, not isolated high logarithmic
jets. The next mixed pairing is what connects those blocks to the product
tangent of the experiment.

\begin{lemma}[Strict confluent mixed pairing]\label{lem:confluent-positive}
Let $\alpha_1<\cdots<\alpha_s$ be nonnegative, and let $\lambda$ run
through distinct nonnegative numbers with multiplicities $q_\lambda$
summing to $s$. Include the columns
$t^\lambda(\log t)^k/k!$, $0\le k<q_\lambda$, in this order within
each increasing $\lambda$ cluster. Assume $q_0=1$ if zero occurs.
The matrix of their integrals against $t^{\alpha_i}\mu(\dd t)$ has
strictly positive determinant. Consequently the corresponding rectangular
pairing has full column rank whenever there are at least as many distinct
row exponents as columns.
\end{lemma}
\begin{proof}
Put $t=e^x$. The functions become $e^{\lambda x}x^k/k!$.
A nonzero linear combination with total multiplicity $s$ has at most
$s-1$ distinct zeros. Indeed divide by the first exponential and
differentiate $q_{\lambda_1}$ times. Rolle's theorem reduces a hypothetical
$s$ zeros to at least $s-q_{\lambda_1}$ zeros. The remaining terms are
exponential polynomials with distinct nonzero exponents and the same
remaining polynomial degrees: the operator $(D+\lambda-\lambda_1)^{q_{\lambda_1}}$
is invertible on polynomials of each fixed degree bound. Induction gives
a contradiction. If no remaining term occurs, the original function is
an exponential times a polynomial of degree less than $s$, and the
assertion is immediate.

Thus the evaluation determinant at distinct ordered points never
vanishes. Its sign is positive. One can determine the sign by letting
ordered simple exponents approach the prescribed repeated exponents and
dividing by their positive within-cluster Vandermonde factors. Equivalently,
the Wronskian of the displayed exponential-polynomial basis has sign
\[
 \prod_{\lambda<\lambda'}(\lambda'-\lambda)^{q_\lambda q_{\lambda'}}>0;
\]
the exponential factor and the factorial normalizations are positive.
Nonvanishing fixes this sign throughout the connected ordered-point domain.

Use the determinant integration identity from
Lemma~\ref{lem:mixed-moment}, with one determinant replaced by this
confluent evaluation determinant. Both determinants have the same positive
sign on ordered distinct positive points. The integrand is nonnegative
elsewhere by symmetry and continuity. Separated ordered subintervals have
positive prior mass, so the integral is strictly positive. The logarithmic
functions are bounded at zero under the stated hypotheses, which also
justifies Fubini. Select any $s$ rows for the rectangular assertion.
\end{proof}

\begin{proposition}[Ordinary Vandermonde exterior scales]\label{prop:ordinary-spectrum}
Fix $q\ge1$ and a multiset $x=(x_1,\ldots,x_q)\in[0,1]^q$.
Let $W(x)=(x_i^{j-1})_{i,j=1}^q$, with singular values
$\sigma_1\ge\cdots\ge\sigma_q\ge0$, and let
\[
 V_\ell(x)=\max_{|J|=\ell}\prod_{\{i,j\}\subset J}|x_i-x_j|,
 \qquad V_1=1.
\]
For $1\le\ell\le q$, uniformly including repeated nodes,
\begin{equation}\label{eq:ordinary-exterior-spectrum}
 V_\ell(x)\asymp_q\prod_{j=1}^{\ell}\sigma_j(W(x))
                =\|\bigwedge^\ell W(x)\|.
\end{equation}
Both sides vanish above the number of distinct sites. This is a
comparison up to dimension-dependent constants, not an equality between
a maximal minor and an exterior norm.
\end{proposition}
\begin{proof}
Permute rows into finite Leja order: choose a first site, then repeatedly
maximize the product of distances to the previously chosen sites. Put
$P_{j-1}(z)=\prod_{i<j}(z-x_i)$ and $d_j=|P_{j-1}(x_j)|$,
with $d_1=1$. Every distance is at most one, so the greedy maximization
and removal of candidates give $d_{j+1}\le d_j$. If there are $s$
distinct sites, the first $s$ pivots are positive and all later pivots
are zero.

Let $T$ have as its $j$th column the coefficients of $P_{j-1}$ in the
power basis. It is upper triangular with diagonal one. Its entries are
bounded elementary symmetric polynomials in numbers from $[0,1]$.
Thus $\|T\|$ is bounded in terms of $q$; writing $T=I+U$ gives
$T^{-1}=\sum_{k=0}^{q-1}(-U)^k$ and the same type of inverse bound.
Evaluation gives $WT=N$, $N_{ij}=P_{j-1}(x_i)$.
For $j\le s$ put $L_{ij}=P_{j-1}(x_i)/d_j$. Greedy maximality
bounds its entries in absolute value by one, and its leading $s$ by $s$
block is lower triangular with diagonal entries $\pm1$. For $j>s$
set the $j$th column of $\widetilde L$ to the coordinate vector $e_j$,
retaining the first $s$ columns of $L$. The resulting block lower
triangular matrix and its inverse have norms bounded solely by $q$.
The corresponding columns of $N$ vanish, since their polynomials already
vanish at every distinct site. Therefore
\[
 W=\widetilde L\operatorname{diag}(d_1,\ldots,d_q)T^{-1}.
\]
Singular-value inequalities for left and right multiplication by an
invertible matrix give $c_qd_j\le\sigma_j(W)\le C_qd_j$ for every
$j$, including the zero tail.

Evaluation of the first $\ell$ monic Newton polynomials on any
$\ell$ sites has determinant equal to their Vandermonde determinant.
It is the selected-row submatrix of $L_{:,1:\ell}$ times
$\operatorname{diag}(d_1,\ldots,d_\ell)$ when $\ell\le s$.
The Leibniz formula bounds its first determinant by $\ell!$, while the
first $\ell$ selected sites give absolute determinant one. Hence
$\prod_{j\le\ell}d_j\le V_\ell\le\ell!\prod_{j\le\ell}d_j$.
For $\ell>s$ both quantities vanish. Combining the comparisons proves
\eqref{eq:ordinary-exterior-spectrum}; the equality with the exterior
norm follows from the singular-value decomposition.
\end{proof}

The preceding proof is the finite comparison in \cite{A1v9note}.
It uses no reciprocal gap and no hypothesis from a clustered Fourier
matrix theorem. Applied to the formal future nodes it identifies the
coordinate content of the scalar determinant profile. The attained
subprobability cube, the global dimension-truncated cover, and the causal
raw update are additional statements proved below; none follows merely
from \eqref{eq:ordinary-exterior-spectrum}.
